% Figure from MATH 590 notes, section 26. Compile with the notes preamble.
\begin{tikzpicture}[scale=1.05, >={Stealth[length=4pt]}]
  % left: onto the figure eight
  \fill[figB!6] (-2.1,-1.5) rectangle (2.1,1.5);
  \foreach \c in {-0.75,0.75}
    \foreach \a in {30,90,150,210,330}
      \draw[figR!80, thick, ->] ($(\c,0)+(\a:0.2)$) -- ($(\c,0)+(\a:0.64)$);
  \foreach \a in {180,135,225}
    \draw[figR!80, thick, ->] ($(-0.75,0)+(\a:1.3)$) -- ($(-0.75,0)+(\a:0.82)$);
  \foreach \a in {0,45,315}
    \draw[figR!80, thick, ->] ($(0.75,0)+(\a:1.3)$) -- ($(0.75,0)+(\a:0.82)$);
  \draw[figR!80, thick, ->] (0,1.35) -- (0,0.4);
  \draw[figR!80, thick, ->] (0,-1.35) -- (0,-0.4);
  \draw[figB, thick] (-0.75,0) circle (0.75) (0.75,0) circle (0.75);
  \fill[figB] (0,0) circle (1.2pt);
  \foreach \x/\l in {-0.75/p, 0.75/q} {
    \filldraw[fill=white, draw=black, thin] (\x,0) circle (1.5pt);
    \node[font=\scriptsize] at (\x,-0.25) {$\l$};
  }
  \node[font=\scriptsize, figB] at (0,-1.7) {figure eight $S^1 \vee S^1$};
  % right: onto the theta space
  \begin{scope}[xshift=4.6cm]
    \fill[figB!6] (-2.1,-1.5) rectangle (2.1,1.5);
    \foreach \a in {20,160,200,340}
      \draw[figR!80, thick, ->] (\a:1.85) -- (\a:1.22);
    \draw[figR!80, thick, ->] (0,1.42) -- (0,1.2);
    \draw[figR!80, thick, ->] (0,-1.42) -- (0,-1.2);
    \foreach \a in {90,150,210,330,30}
      \draw[figR!80, thick, ->] ($(-0.55,0)+(\a:0.2)$) -- ($(-0.55,0)+(\a:0.44)$);
    \foreach \a in {90,150,210,330,30}
      \draw[figR!80, thick, ->] ($(0.55,0)+(\a:0.2)$) -- ($(0.55,0)+(\a:0.44)$);
    \draw[figB, thick] (0,0) circle (1.1);
    \draw[figB, thick] (0,1.1) -- (0,-1.1);
    \fill[figB] (0,1.1) circle (1.2pt) (0,-1.1) circle (1.2pt);
    \foreach \x/\l in {-0.55/p, 0.55/q} {
      \filldraw[fill=white, draw=black, thin] (\x,0) circle (1.5pt);
      \node[font=\scriptsize] at (\x,-0.25) {$\l$};
    }
    \node[font=\scriptsize, figB] at (0,-1.7) {theta space $\Theta$};
  \end{scope}
\end{tikzpicture}
